% Opening results and conventions, to follow V4's Motivation slide.
% Numerical provenance and qualifications: SOURCES.md and review/.
% Generated by questions_audit/make_slide4_no_cut.py: no GNN score cut; reco >=2 tracks.
\newcommand{\gnnEfficiency}{$69.06\pm0.06\%$}
\newcommand{\gnnPurity}{$69.69\pm0.06\%$}
\newcommand{\gnnAssigned}{$99.10\pm0.01\%$}
\newcommand{\gnnCleanReco}{$49.75\pm0.13\%$}
\newcommand{\amvfEfficiency}{$68.32\pm0.06\%$}
\newcommand{\amvfPurity}{$68.65\pm0.06\%$}
\newcommand{\amvfAssigned}{$99.52\pm0.01\%$}
\newcommand{\amvfCleanReco}{$35.17\pm0.10\%$}
\newcommand{\gnnGainEfficiency}{$+0.75\pm0.02$}
\newcommand{\gnnGainPurity}{$+1.05\pm0.02$}

\newcommand{\sourcecaption}[1]{\par\vspace{.5em}{\scriptsize\color{black!55}#1\par}}
\newcommand{\itkpaper}{\href{https://arxiv.org/html/2412.15090v2\#S5.SS5}{ITk performance paper, Sec.\ 5.5}}
\newcommand{\runthreepaper}{\href{https://arxiv.org/abs/2308.09471}{Run 3 tracking paper, Fig.\ 13}}
\newcommand{\amvfnote}{\href{https://cds.cern.ch/record/2670380}{AMVF note (2019), Sec.\ 3.2}}

\begin{frame}{PV-Finder and AMVF: vertex finding}
\context{$14$ TeV $t\bar t$ simulation, $\langle\mu\rangle=192.5$, $|\eta|<4$; same 1,920 events; integral cut $0.3$.}
\centering
\small\renewcommand{\arraystretch}{1.2}
\begin{tabular}{lccc}
\toprule
Measurement & $z$ window & PV-Finder & AMVF\\
\midrule
All-vertex efficiency & $\pm0.1$ mm & $74.11\pm0.10\%$ & $70.99\pm0.10\%$\\
Unmatched reco / event & $\pm0.1$ mm & $25.16\pm0.11$ & $25.46\pm0.11$\\
Extra nearby reco / event & $\pm0.1$ mm & $0.112\pm0.008$ & $0.035\pm0.004$\\
\midrule
Hard-scatter efficiency & $\pm0.1$ mm & $99.90^{+0.07}_{-0.14}\%$ & $99.79^{+0.10}_{-0.16}\%$\\
\bottomrule
\end{tabular}
\par\vspace{.25em}\raggedright\footnotesize
\textbf{Efficiency:} truth PVs with $\ntrk\geq2$ and a reco vertex within $0.1$ mm / all truth PVs with $\ntrk\geq2$. One reco vertex can be near several truth PVs.\\[.3em]
\textbf{Unmatched:} no eligible truth PV within $0.1$ mm. \textbf{Extra nearby reco:} another candidate near an already matched truth PV.\\[.3em]
\textbf{Hard scatter:} events with a vertex within $0.1$ mm of the true HS position / all events.
\sourcecaption{PV-Finder: peaks before GNN; AMVF: $\ntrk\geq2$. Unmatched candidates may lie near excluded one-track truth PVs. Errors: event bootstrap (10,000 resamples); HS: 68.27\% Clopper--Pearson.}
\end{frame}

\begin{frame}{One set of settings}
\vspace{.8em}
\centering
{\Large Integral threshold: \textbf{\textcolor{pennred}{0.3}}}\par
\vspace{1.8em}
{\Large Vertex finding: \textbf{\textcolor{pennred}{$|\Delta z|\leq0.1$ mm}}}\par
\vspace{1.8em}
{\Large Vertex categories: \textbf{\textcolor{pennred}{ATLAS convention}}}\par
\vspace{.45em}
{\small Clean / Matched, Merged, Split, Fake}\par
\vfill
\sourcecaption{0.5 is a cross-check only. Current category counts approximate ATLAS: unit track weights and an unlabelled-track proxy for Fake.}
\end{frame}

\begin{frame}{HL-LHC: resolution at integral cuts 0.3 and 0.5}
\centering
\includegraphics[width=.99\linewidth,height=.78\textheight,keepaspectratio]{figures/resolution_integral_030_050.pdf}
\sourcecaption{AMVF is unchanged between panels. This reconstructed-pair fit uses no truth-matching window.}
\end{frame}

\begin{frame}{One vertex classification: the ATLAS convention}
\context{Same four categories as V4 slide 26; truth contributions use vertex-fit weights.}
\small\renewcommand{\arraystretch}{1.5}
\begin{tabular}{p{3.2cm}p{10.3cm}}
\toprule
\textbf{Clean / Matched} & At least 70\% of the total track weight comes from one true collision.\\
\textbf{Merged} & No true collision contributes 70\% of the total track weight.\\
\textbf{Split} & Several reco vertices have the same dominant true collision. Only the vertex with the largest $\sum p_T^2$ keeps Clean or Merged.\\
\textbf{Fake} & Fake tracks contribute more weight than any single true collision.\\
\bottomrule
\end{tabular}
\par\vspace{.6em}\footnotesize
Fake and Split take precedence over the Clean/Merged label. We retain reco vertices with at least two tracks for both methods.
\sourcecaption{\amvfnote. Current plots use unit weights and unlabelled tracks as a Fake proxy; the ntuples lack fit weights and detailed fake-track truth information. Position-only results use ``nearby'' and ``unmatched'', not these category names.}
\end{frame}


\begin{frame}{PV-Finder + GNN: track association}
\context{Same 1,920 events, $185\leq\mu\leq200$; integral cut $0.3$; no GNN score cut; reco $\ntrk\geq2$ for both.}
\centering\small\renewcommand{\arraystretch}{1.2}
\begin{tabular}{lcc}
\toprule
Measurement & PV-Finder + GNN & AMVF (stored links)\\
\midrule
Track efficiency & \gnnEfficiency & \amvfEfficiency\\
Assignment purity & \gnnPurity & \amvfPurity\\
Assigned fraction & \gnnAssigned & \amvfAssigned\\
\midrule
Clean / reconstructed vertices & \gnnCleanReco & \amvfCleanReco\\
\bottomrule
\end{tabular}
\par\vspace{.25em}\raggedright\footnotesize
\textbf{Truth-labelled tracks:} efficiency = correct / all; purity = correct / assigned links. All stored truth origins enter these track metrics.\\[.3em]
Correct means the vertex's dominant truth origin matches the track's origin.\\[.3em]
Clean uses 70\% purity from track counts; the public definition uses fit weights.
\takeaway{At high pile-up: \gnnGainEfficiency\ efficiency points and \gnnGainPurity\ purity points over AMVF.}
\sourcecaption{Highest-scoring target per track; then remove reco vertices with fewer than two tracks. Paired event-bootstrap errors. Categories approximate \amvfnote; fit weights are unavailable.}
\end{frame}

\begin{frame}{Proposed paper structure}
\begin{columns}[T]
\begin{column}{.45\linewidth}
\small
\textbf{1\quad Introduction}\\[.45em]
Why vertex reconstruction matters at high pile-up.\\[1em]
\textbf{2\quad ATLAS detector}\\[.6em]
\textbf{3\quad Run 3 Inner Detector}\\[.6em]
\textbf{4\quad The ITk upgrade}\\[1em]
\textbf{5\quad Data and MC samples}\\[.45em]
Run 3 data and simulation; HL-LHC simulation.\\[.35em]
Track selection, pile-up and beam spot.
\end{column}
\begin{column}{.51\linewidth}
\small
\textbf{6\quad PV-Finder and the track-to-vertex GNN}\\[.45em]
Inputs, both architectures, training and chosen settings.\\[1em]
\textbf{7\quad Performance}\\[.45em]
7.1\quad Run 3 simulation\\[.3em]
7.2\quad Run 3 data\\[.3em]
7.3\quad HL-LHC simulation\\[.3em]
7.4\quad Timing studies\\[1em]
\textbf{8\quad Conclusion}\\[.45em]
Main improvements and remaining limitations.
\end{column}
\end{columns}
\takeaway{Describe both algorithms in Section 6; compare their performance with AMVF in Section 7.}
\sourcecaption{Section order follows the professor's outline, dated 23 September 2026.}
\end{frame}

\begin{frame}{What we will include in the paper}
\small
\textbf{Use an integral cut of 0.3 for the main results.}\\[.4em]
Use the same setting for vertex yields, resolution and the GNN input vertices.
\par\vspace{.6em}
\centering\footnotesize\renewcommand{\arraystretch}{1.2}
\begin{tabular}{p{2.1cm}p{4.35cm}p{6.35cm}}
\toprule
Study & Available & Add or update\\
\midrule
Run 3 MC & V4 has Run 2 MC validation. & Add representative Run 3 MC; finder and GNN against AMVF.\\
Run 3 data & Agreement with AMVF positions and track assignments. & Compare the same observables in data and MC, with matched pile-up.\\
HL-LHC MC & Inclusive and HS efficiency; TTVA efficiency, purity and categories. & Add HS selection and HS position resolution; check the AMVF fit weights.\\
Timing & Athena CPU and GPU measurements. & Compare CPU with CPU; show finder, GNN and full-chain times versus pile-up.\\
\bottomrule
\end{tabular}
\par\vspace{.5em}\raggedright\small
Use one set of vertex categories throughout: the ATLAS convention.
\sourcecaption{Main GNN results use no score cut; the score scan is a separate diagnostic. Re-evaluate Run 2/3 at integral cut 0.3 before quoting updated performance.}
\end{frame}

\begin{frame}{Comparison with the ITk paper}
\centering\small\renewcommand{\arraystretch}{1.25}
\begin{tabular}{p{3.05cm}p{3.05cm}p{6.6cm}}
\toprule
Measurement & Published result & Comparison with our study\\
\midrule
HS finding, $0.1$ mm & About $99.5\%$ & PV-Finder $99.90\%$, AMVF $99.79\%$. Same $z$ cut; exact monitoring selection to check.\\
HS finding + selection & About $92$--$93\%$ & Stored AMVF primary: $93.0\%$. The old $80.7\%$ re-ranked the ntuple track lists.\\
HS position precision & ITk: about $10\,\mu$m; Run 3: about $20\,\mu$m & Add HS resolution. Our inclusive width and vertex-separation scale are different quantities.\\
\bottomrule
\end{tabular}
\par\vspace{.6em}\raggedright\small
The stored primary contains the largest number of true HS tracks in 1,785/1,920 events ($93.0^{+0.6}_{-0.6}\%$). Re-ranking changes the chosen vertex in 428 events.
\takeaway{The stored AMVF primary gives HS selection consistent with the published range.}
\sourcecaption{\itkpaper, Figs.\ 26--28. Our ntuple uses RecoVertex\_type = 1. Confirm geometry and the exact monitoring configuration before claiming a reproduction.}
\end{frame}

\begin{frame}{Run 3 timing: the published comparison}
\begin{columns}[c]
\begin{column}{.43\linewidth}
\small
\textbf{ACTS AMVF: about 200 ms/event at $\langle\mu\rangle\simeq90$.}\\[.7em]
That point is from a 2018 machine-development run, outside the standard data-quality selection.\\[.7em]
Benchmark: AMD EPYC 7302, one thread per reconstruction job.
\end{column}
\begin{column}{.54\linewidth}
\centering
\includegraphics[page=22,trim={172bp 562bp 157bp 104bp},clip,width=.96\linewidth,height=.48\textheight,keepaspectratio]{research/run3_tracking.pdf}
\end{column}
\end{columns}
\par\vspace{.9em}\small
V4's approximately $9\times$ gain compares \textbf{PV-Finder + GNN on GPU} with \textbf{AMVF on CPU}, on different hardware from the paper (Xeon 4216 / Tesla T4).
\par\vspace{.55em}
For the paper, show a common benchmark and separate finder, GNN and full-chain times.
\sourcecaption{\runthreepaper\ and benchmark description in Sec.\ 4; V4 timing source: \treysource.}
\end{frame}
